% Generated draft: requires mathematical and source review.
\documentclass[11pt]{article}
\usepackage[margin=25mm]{geometry}
\usepackage{fontspec}
\setmainfont{Latin Modern Roman}
\usepackage{amsmath,amssymb,mathtools,unicode-math}
\setmathfont{Latin Modern Math}
\usepackage{enumitem,needspace,xurl,hyperref}
\setlength{\parindent}{0pt}
\setlength{\parskip}{5pt}
\setlength{\emergencystretch}{3em}
\newcommand{\N}{\mathbb{N}}
\newcommand{\Z}{\mathbb{Z}}
\newcommand{\Q}{\mathbb{Q}}
\newcommand{\R}{\mathbb{R}}
\newcommand{\C}{\mathbb{C}}
\newcommand{\F}{\mathbb{F}}
\newcommand{\A}{\mathbb{A}}
\newcommand{\PP}{\mathbb{P}}
\newcommand{\E}{\mathbb{E}}
\newcommand{\eps}{\varepsilon}
\newcommand{\dd}{\,\mathrm{d}}
\begin{document}
\section*{UnsolvedMath ID: 30000599}
\textbf{Source problem code:} catalogue entry
\subsection{Definitions and mathematical statement}
\textbf{Periodic-measure conjecture and construction questions.} For a plane billiard table, the set of periodic billiard trajectories should have measure zero for the canonical invariant measure. Separately, are convex tables with an invariant curve consisting of periodic trajectories rare or common, and how can such tables be constructed systematically? The source asks analogous questions for outer billiards as well. All these requests are retained.\par In a conventional smooth bounded inner-billiard setting, the phase space consists of boundary impact points \(s\) and inward unit directions, with signed angle \(\theta\in(-\pi/2,\pi/2)\) measured from the inward normal. The billiard map \(B\) sends one regular impact to the next by specular reflection and preserves \(d\mu=\cos\theta\,ds\,d\theta\), up to normalization. Singular grazing or other undefined iterates must be treated under the table's regularity assumptions. The conjecture is\[\mu\left(\bigcup_{q\ge1}\{z:B^qz=z\}\right)=0,\] on the regular phase space. This measures initial conditions, not the number of geometric orbit shapes, and is equivalent to measure zero for every fixed-period set. The source does not specify table regularity or boundary singularities, so the displayed smooth convention is conditional. Convexity is explicitly requested for the invariant-curve questions, not silently added to every version of the measure conjecture.\par An invariant periodic curve means an appropriate curve \(C\) in phase space with \(B(C)=C\) and every point of \(C\) periodic. Its smoothness, graph or caustic requirements, and any common-period condition must be supplied. A one-dimensional regular curve in this two-dimensional phase space can have zero invariant area, so its existence does not contradict the first question. The terms \emph{rare} and \emph{common} require a specified topology, genericity notion, or measure on the class of convex tables; none is given. Systematic construction also needs the admissible regularity and geometric constraints.\par For a conventional outer billiard outside a smooth strictly convex planar body, the map sends an exterior point to its reflection in the selected tangent contact point, with a fixed orientation convention, and the relevant invariant measure is planar area. The source does not fix this regularity or orientation, nor distinguish smooth from polygonal bodies. The outer analogues must therefore be formulated in their own phase space and class; the inner conjecture is not automatically a universal outer-billiard assertion.\par The supplied triage checked 2026-08-21 labels the measure conjecture open for arbitrary smooth bounded domains, with special classes settled, and links it to the periodic-orbit hypothesis in two-term Weyl asymptotics. Its wording about equivalence is not accompanied by a precise spectral theorem, so it is not turned into an unsupported if-and-only-if between a spectral expansion and periodic-measure vanishing. The assessment does not separately resolve the invariant-curve or outer questions and requires independent verification.
\subsection{Short English statement}
The record combines a zero-measure conjecture for periodic inner-billiard initial conditions with questions about periodic invariant curves and analogous outer dynamics. Such curves can have zero area. The relevant phase spaces, table regularity, and meaning of rarity must be kept distinct.
\subsection{Sources}
\begin{itemize}\item Local UnsolvedMath registry, permanent ID 30000599, code catalogue entry; \emph{Periodic Trajectories and Invariant Curves in Billiards}, Oberwolfach Reports Open Problems.\item \emph{Geometrie (2007)}, Oberwolfach Reports, DOI 10.4171/owr/2006/45, pp. 2699--2756. \url{https://doi.org/10.4171/owr/2006/45}. Proposed year and year tag: 2006; no proposer named.\item Supplied literature label \emph{original OWR report}: \url{https://doi.org/10.4171/owr/2006/45}.\item Supplied literature label \emph{spectral geometry survey}: \url{https://nsa.fjfi.cvut.cz/david/other/course-Toulouse.pdf}.\item Literature triage checked 2026-08-21, status \texttt{open}; record updated 2026-09-16. The \texttt{corrected\_verified} field records preserving the measure claim, both construction questions, and outer extension. Verify dated claims independently.\end{itemize}
\end{document}
